\chapterbackmatter{Supplementary Exercises}

\begin{enumerate}
  \SupplementaryExercise{1}{Connected Vacuum-Graph Exponentiation}
  {Adapted from Ian Lim, Advanced Quantum Field Theory, Section 5 (2019).}
  {lim-advanced-qft-2019}
  Let \(Z[J]\) be the sum of all vacuum graphs in the presence of a source.
  By decomposing each graph into connected components and accounting for
  permutations of identical components, prove that
  \(Z[J]=\exp(iW[J])\), where \(iW[J]\) is the sum of connected vacuum
  graphs.  Explain why derivatives of \(W\), rather than \(Z\), generate
  connected correlation functions.

  \SupplementaryExercise{2}{Legendre Hessians as Inverse Kernels}
  {Adapted from John McGreevy, UC San Diego Physics 215C lecture notes, Section 2.1 (2014).}
  {mcgreevy-physics215c-2014}
  Starting from
  \[
    \phi^i=\frac{\delta W}{\delta J_i},\qquad
    \Gamma[\phi]=W[J_\phi]-\phi^iJ_{\phi i},
  \]
  derive \(\Gamma_{,ij}W_{,jk}=-\delta_i{}^k\) in collective-index
  notation.  State the assumptions needed to invert the two Hessians and
  interpret each kernel diagrammatically.

  \SupplementaryExercise{3}{A Connected Three-Point Tree}
  {Adapted from Ian Lim, Advanced Quantum Field Theory, Section 5 (2019).}
  {lim-advanced-qft-2019}
  Differentiate the inverse-Hessian identity once with respect to the
  current.  Express \(W_{,ijk}\) in terms of the full propagator and the
  1PI three-point vertex, keeping the sign convention
  \(W^{(2)}=-\Gamma^{(2)-1}\).  Draw or describe the corresponding tree
  and include its external propagators.

  \SupplementaryExercise{4}{Four-Point Channels from the Quantum Action}
  {Adapted from Mark Srednicki, Quantum Field Theory prepublication draft, Sections 17 and 21 (2006).}
  {srednicki-qft-draft-2006}
  Obtain the complete relation between \(W_{,ijkl}\) and derivatives of
  \(\Gamma\).  Show that it contains one 1PI four-point term and three
  exchange terms made from pairs of 1PI three-point vertices.  For four
  identical scalars, identify the three inequivalent partitions of the
  external labels.

  \SupplementaryExercise{5}{Loop Counting with an Auxiliary Parameter}
  {Adapted from Howard Haber, Physics 222, Functional Integral Derivation of the Effective Potential, pp. 1--4 (2020).}
  {ucsc-physics222-effective-potential-2020}
  Replace the action in a functional integral by \(g^{-1}\Gamma[\phi]\).
  Show that a graph with \(I\) internal lines and \(V\) vertices carries
  \(g^{I-V}\).  Use \(L=I-V+1\) for a connected graph to derive the
  \(g^{L-1}\) loop counting of \(W_\Gamma[J,g]\).  Explain why the
  \(g\to0\) saddle retains exactly the connected trees.

  \SupplementaryExercise{6}{Zero-Momentum Vertices and the Potential}
  {Adapted from Howard Haber, Physics 222, Functional Integral Derivation of the Effective Potential, p. 1 (2020).}
  {ucsc-physics222-effective-potential-2020}
  For a translation-invariant scalar theory, expand \(\Gamma[\phi]\) in
  1PI \(n\)-point kernels about \(\phi=0\).  Set \(\phi(x)=\phi_0\) and
  prove that the effective potential is determined by the zero-momentum
  vertices.  Carefully track the spacetime-volume delta function and the
  sign in \(\Gamma[\phi_0]=-\mathcal V_4V(\phi_0)\).

  \SupplementaryExercise{7}{The Scalar Gaussian Determinant}
  {Adapted from Roman Jackiw, Functional Evaluation of the Effective Potential, Sections II--III (1974).}
  {jackiw-effective-potential-1974}
  Expand a real scalar action about a constant background and retain the
  quadratic fluctuation operator \(K\).  Evaluate the Gaussian integral to
  show that the one-loop effective action is proportional to
  \(\frac{i}{2}\operatorname{Tr}\ln K\).  After Wick rotation, differentiate
  the four-dimensional momentum integral three times with respect to the
  field-dependent mass squared and recover the logarithmic term.

  \SupplementaryExercise{8}{The Fermion-Loop Minus Sign}
  {Adapted from Hong Liu, MIT OpenCourseWare 8.323, Problem Set 10, Problem 3 (2023).}
  {mit-8323-s23-ps10}
  Evaluate a finite-dimensional Grassmann Gaussian and then extend the
  result formally to a Dirac field.  Explain why a complex fermion produces
  \(\det D\), whereas a real boson produces \((\det K)^{-1/2}\).  Reduce
  \(\operatorname{Tr}\ln(\slashed p-M)\) to a scalar logarithm and identify
  the spin multiplicity.

  \SupplementaryExercise{9}{One-Loop Supertrace Formula}
  {Adapted from Howard Haber, Physics 222, Functional Integral Derivation of the Effective Potential, pp. 8--10 (2020).}
  {ucsc-physics222-effective-potential-2020}
  Consider constant backgrounds in a theory whose bosonic and fermionic
  quadratic fluctuation operators have mass-squared eigenvalues
  \(m_b^2(\phi)\) and \(m_f^2(\phi)\).  Derive a compact one-loop expression
  using a graded trace over species.  State how real scalars, complex
  scalars, and Dirac fermions are weighted, and explain which polynomial
  terms are scheme dependent.

  \SupplementaryExercise{10}{Radial and Goldstone Fluctuation Masses}
  {Adapted from Roman Jackiw, Functional Evaluation of the Effective Potential, Section III (1974).}
  {jackiw-effective-potential-1974}
  For
  \[
    V_0(\boldsymbol\phi)
      =\frac12m^2\boldsymbol\phi^2
       +\frac{\lambda}{4!}(\boldsymbol\phi^2)^2
  \]
  in an \(O(N)\) theory, choose the background
  \(\boldsymbol\phi=(\varphi,0,\ldots,0)\).  Diagonalize the Hessian and
  find the radial and \(N-1\) transverse mass-squared eigenvalues.  Use
  them to write the nonpolynomial part of the four-dimensional one-loop
  potential.

  \SupplementaryExercise{11}{Complex-Scalar Background Determinant}
  {Adapted from Howard Haber, Physics 222, Functional Integral Derivation of the Effective Potential, pp. 8--11 (2020).}
  {ucsc-physics222-effective-potential-2020}
  Let a complex scalar have
  \(V_0=m^2\Phi^*\Phi+\lambda(\Phi^*\Phi)^2/4\).  Choose a real constant
  background and rewrite the two real fluctuations as radial and angular
  modes.  Determine their field-dependent masses and the one-loop
  logarithms.  Verify explicitly that the answer depends only on
  \(\Phi^*\Phi\).

  \SupplementaryExercise{12}{A Yukawa Background Determinant}
  {Adapted from Hong Liu, MIT OpenCourseWare 8.323, Problem Set 10, Problem 4 (2023).}
  {mit-8323-s23-ps10}
  In four dimensions take a real scalar background \(\phi\) and a Dirac
  fermion with mass \(M(\phi)=m+y\phi\).  Compute the nonpolynomial
  fermion-loop contribution to \(V_{\rm eff}(\phi)\).  List all local
  counterterms required when no reflection symmetry is imposed, and show
  how the list is reduced if a symmetry forces \(M(-\phi)=-M(\phi)\).

  \SupplementaryExercise{13}{Renormalization Conditions for \(\phi^4\)}
  {Adapted from Vadim Kaplunovsky, Vacuum Energy and Effective Potentials, pp. 3--8 (2019).}
  {kaplunovsky-effective-potential-2019}
  For the one-loop potential of a real \(\phi^4\) theory, impose
  \[
    V(0)=0,\qquad V''(0)=m_R^2,\qquad
    V^{(4)}(\phi)\big|_{\phi=\kappa}=g_R .
  \]
  Show how these conditions determine the three polynomial integration
  constants left by differentiating the logarithmic integral.  Explain why
  changing the subtraction point changes polynomial terms but not the
  unrenormalized determinant.

  \SupplementaryExercise{14}{Dimensional Transmutation at One Loop}
  {Adapted from Sidney Coleman and Erick Weinberg, Radiative Corrections as the Origin of Spontaneous Symmetry Breaking, Sections II--III (1973).}
  {coleman-weinberg-1973}
  Consider a massless theory whose one-loop potential can be written
  \[
    V(\phi)=A\phi^4+B\phi^4\ln(\phi^2/\kappa^2),
    \qquad B>0.
  \]
  Find every nonzero stationary point and determine when it is a local
  minimum.  Eliminate \(A\) in favor of the nonzero minimum \(v\), rewrite
  the potential using \(v\), and explain how a dimensionless coupling has
  been traded for a mass scale.

  \SupplementaryExercise{15}{The Logarithm in General Even Dimension}
  {Adapted from Mark Srednicki, Quantum Field Theory prepublication draft, Section 19 (2006).}
  {srednicki-qft-draft-2006}
  Define
  \[
    \mathcal I_{2n}(x)
      =\frac12\int\frac{d^{2n}k_E}{(2\pi)^{2n}}\ln(k_E^2+x).
  \]
  Differentiate sufficiently many times to obtain a convergent integral
  and show that its nonpolynomial part is proportional to
  \((-1)^n x^n\ln(x/\kappa^2)\).  Determine the coefficient and check the
  cases \(D=4\) and \(D=6\).

  \SupplementaryExercise{16}{Which Counterterms Does the Determinant Demand?}
  {Adapted from Roman Jackiw, Functional Evaluation of the Effective Potential, Section II (1974).}
  {jackiw-effective-potential-1974}
  In \(D\) spacetime dimensions suppose the fluctuation mass is affine,
  \(\mu^2(\phi)=a+b\phi\).  Use power counting of
  \(\frac12\operatorname{Tr}\ln(-\partial^2+\mu^2)\) to determine the
  degree of the divergent polynomial in \(\mu^2\) for even \(D\).
  Specialize to \(D=6\) and explain why vacuum, tadpole, mass, and cubic
  counterterms form a closed set.

  \SupplementaryExercise{17}{Imaginary Part in a Spinodal Region}
  {Adapted from John McGreevy, UC San Diego Physics 215C lecture notes, Section 2.2 (2014).}
  {mcgreevy-physics215c-2014}
  Analytically continue the one-loop scalar term
  \(m^4(\phi)\ln[m^2(\phi)-i0]/(64\pi^2)\) into a region where
  \(m^2(\phi)<0\).  Compute its imaginary part and relate its sign to the
  chosen \(i0\) prescription.  Explain physically why this perturbative
  result does not contradict the reality of the exact constrained energy.

  \SupplementaryExercise{18}{Convexity from the Legendre Transform}
  {Adapted from Ian Lim, Advanced Quantum Field Theory, Section 5 (2019).}
  {lim-advanced-qft-2019}
  Work in Euclidean signature, where the connected two-point kernel is
  positive semidefinite.  Use the inverse-Hessian relation and
  \(\Gamma[\phi_0]=-\mathcal V_4V(\phi_0)\) to show that the exact
  effective potential is convex.  Discuss zero modes and state why a
  finite-order loop expansion need not display this property.

  \SupplementaryExercise{19}{The Maxwell Segment of the Exact Potential}
  {Adapted from Vadim Kaplunovsky, Vacuum Energy and Effective Potentials, pp. 8--11 (2019).}
  {kaplunovsky-effective-potential-2019}
  Suppose two homogeneous phases have order parameters \(\phi_1<\phi_2\)
  and energy densities \(V_1,V_2\).  Construct a large-volume mixed phase
  with prescribed average
  \(\phi=\alpha\phi_1+(1-\alpha)\phi_2\).  Neglecting interface energy per
  unit volume, derive the linear segment connecting the two endpoints.
  Explain why it is horizontal for degenerate vacua.

  \SupplementaryExercise{20}{Static Sources and Constrained Energy}
  {Adapted from John McGreevy, UC San Diego Physics 215C lecture notes, Section 2.2 (2014).}
  {mcgreevy-physics215c-2014}
  Turn on a time-independent source adiabatically for a long interval
  \(T\).  Show that \(W[J]=-E[J]T\) up to switching effects.  Legendre
  transform with respect to the static source and prove that
  \(-\Gamma[\phi]/T\) is the minimum energy among normalized states with
  the prescribed field expectation value.  State where the adiabatic and
  infinite-volume assumptions enter.

  \SupplementaryExercise{21}{First Terms in the Derivative Expansion}
  {Adapted from Howard Haber, Physics 222, Functional Integral Derivation of the Effective Potential, pp. 1--3 (2020).}
  {ucsc-physics222-effective-potential-2020}
  For a slowly varying real scalar background, write the most general
  Lorentz-invariant effective action through two derivatives as
  \[
    \Gamma[\phi]=\int d^4x\,
      \left[-V(\phi)-\frac12 Z(\phi)
      \partial_\mu\phi\,\partial^\mu\phi\right]+O(\partial^4).
  \]
  Explain which momentum derivatives of the 1PI two-point vertex determine
  \(Z(\phi)\), why a constant background determines only \(V\), and how a
  local field redefinition changes \(Z\) and \(V\).

  \SupplementaryExercise{22}{Finite Linear Symmetry of the Quantum Action}
  {Adapted from Ian Lim, Advanced Quantum Field Theory, Section 5 (2019).}
  {lim-advanced-qft-2019}
  Let a finite invertible matrix \(M\) leave the action and functional
  measure invariant under \(\phi\mapsto M\phi\).  Derive the induced
  transformation of \(J\), prove the corresponding identity for \(W[J]\),
  and then prove \(\Gamma[M\phi]=\Gamma[\phi]\).  Do not assume that \(M\)
  is orthogonal.

  \SupplementaryExercise{23}{The \(O(N)\) Ward Identity for \(\Gamma\)}
  {Adapted from Mark Srednicki, Quantum Field Theory prepublication draft, Sections 22 and 24 (2006).}
  {srednicki-qft-draft-2006}
  For infinitesimal \(O(N)\) rotations
  \(\delta\phi_a=\epsilon^A(T^A)_{ab}\phi_b\), use a change of variables in
  the source functional to derive
  \[
    \int d^4x\,\frac{\delta\Gamma}{\delta\phi_a(x)}
      (T^A)_{ab}\phi_b(x)=0.
  \]
  Differentiate this identity once and twice with respect to the field and
  state the resulting constraints on 1PI vertices at a symmetric
  background.

  \SupplementaryExercise{24}{Goldstone Zero Modes from \(\Gamma^{(2)}\)}
  {Adapted from John McGreevy, UC San Diego Physics 215C lecture notes, Section 2.3 (2014).}
  {mcgreevy-physics215c-2014}
  Let a constant stationary point \(v_a\) of the exact effective action
  break a continuous linear symmetry, so that
  \(T^Av\ne0\).  Differentiate the Ward identity and prove that
  \(\Gamma^{(2)}_{ab}(p=0)(T^Av)_b=0\).  Interpret this zero eigenvalue as a
  massless excitation and count the independent modes in terms of the
  broken generators.

  \SupplementaryExercise{25}{Reflection Symmetry and Odd 1PI Vertices}
  {Adapted from Ian Lim, Advanced Quantum Field Theory, Section 5 (2019).}
  {lim-advanced-qft-2019}
  Suppose the action and measure are invariant under
  \(\phi\mapsto-\phi\).  Prove successively that \(Z[J]=Z[-J]\),
  \(\phi_{-J}=-\phi_J\), and \(\Gamma[-\phi]=\Gamma[\phi]\).
  Show that every odd 1PI vertex at the symmetric configuration vanishes,
  and state how the conclusion changes after expansion about a nonzero
  symmetry-breaking minimum.

  \SupplementaryExercise{26}{Nonlinear Symmetry and an Averaging Obstruction}
  {Adapted from John McGreevy, UC San Diego Physics 215C lecture notes, Sections 2.1--2.3 (2014).}
  {mcgreevy-physics215c-2014}
  Let a measure-preserving infinitesimal change of variables be
  \(\delta\chi^n=\epsilon F^n[\chi]\).  Derive the source Ward identity and
  show that the effective action is invariant under
  \(\delta\chi^n=\epsilon\langle F^n[\widehat\chi]\rangle_{J_\chi}\).
  For \(F[\chi]=\chi^2\), express the difference between
  \(\langle F[\widehat\chi]\rangle\) and \(F[\chi]\) in terms of the
  connected two-point function.

  \SupplementaryExercise{27}{A Jacobian Term in a Ward Identity}
  {Adapted from John McGreevy, UC San Diego Physics 215C lecture notes, Section 5.1 (2014).}
  {mcgreevy-physics215c-2014}
  Repeat the infinitesimal change-of-variables derivation when the
  functional measure has Jacobian
  \(\mathcal D\chi'\!=\mathcal D\chi\,
  [1+i\epsilon\mathcal A[\chi]]\).  Derive the modified Ward identity for
  \(W\) and its effective-action form.  Explain why a nonremovable local
  functional \(\mathcal A\) is called an anomaly and why adding an allowed
  local counterterm can sometimes shift it.

  \SupplementaryExercise{28}{Graded Legendre Transform}
  {Adapted from Hong Liu, MIT OpenCourseWare 8.323, Problem Set 10, Problem 3 (2023).}
  {mit-8323-s23-ps10}
  Introduce Grassmann sources for fermionic fields, with every current
  written to the right of its field.  Using explicit left and right
  derivatives, derive
  \[
    \frac{\delta_R W}{\delta J_n}=\chi^n,\qquad
    \frac{\delta_L\Gamma}{\delta\chi^n}=-J_n.
  \]
  Work out the graded inverse-Hessian relation for one bosonic and one
  fermionic species and identify where ordinary matrix transposition must
  be replaced by graded transposition.

  \SupplementaryExercise{29}{A Source for a Composite Operator}
  {Adapted from John Cornwall, Roman Jackiw, and E. Tomboulis, Effective Action for Composite Operators (1974).}
  {cornwall-jackiw-tomboulis-1974}
  Add a source \(K(x)\) for the local composite operator
  \(O(x)=\phi^2(x)/2\), and Legendre transform \(W[J,K]\) with respect to
  both sources.  Derive the two stationarity equations.  Explain why the
  resulting functional cannot be identified with the ordinary elementary
  field 1PI action alone, and determine which additional local
  counterterms involving \(K\) are allowed by power counting in four
  dimensions.

  \SupplementaryExercise{30}{Scale Independence of the Effective Potential}
  {Adapted from Sidney Coleman and Erick Weinberg, Radiative Corrections as the Origin of Spontaneous Symmetry Breaking, Sections III--IV (1973).}
  {coleman-weinberg-1973}
  Starting from the fact that the bare effective potential is independent
  of the subtraction scale, derive
  \[
    \left(\kappa\frac{\partial}{\partial\kappa}
      +\sum_i\beta_i\frac{\partial}{\partial g_i}
      -\gamma\,\phi\frac{\partial}{\partial\phi}\right)V(\phi)=0.
  \]
  Apply it to
  \(V=\lambda\phi^4/4!+B\phi^4\ln(\phi^2/\kappa^2)\) through one loop and
  find the relation among \(B\), the one-loop beta function, and the
  anomalous dimension.
\end{enumerate}
